\section*{Аннотация}

В выпускной квалификационной работе рассматривается задача траекторного управления нелинейным динамическим объектом в трехмерном пространстве. В качестве объекта исследования выбран квадрокоптер, для которого построена программная модель движения, реализован базовый алгоритм согласованного управления по выходу и разработан интеллектуальный модуль подбора параметрической скорости движения вдоль заданной кривой.

Основной акцент сделан на связке двух уровней управления. Нижний уровень отвечает за устойчивое слежение за пространственной траекторией и формирует управляющие воздействия с учетом нелинейной динамики объекта. Верхний уровень подбирает значение параметрической скорости $V^{\ast}$ в зависимости от текущего состояния системы и локальной геометрии траектории. Для этого был сформирован датасет на основе оракульного поиска допустимой скорости и обучена полносвязная нейронная сеть.

В работе описаны математическая модель объекта, структура регулятора, состав признаков для интеллектуального модуля, порядок формирования датасета и результаты численного моделирования. Проведенные эксперименты показывают, что предложенный подход позволяет увеличить среднюю скорость движения на тестовой спиральной траектории без потери устойчивости и с сохранением приемлемой точности слежения.

\newpage
